% results.tex (국문 v11). 차이 = 앞 방법의 RMSE − 뒤 방법의 RMSE(같은 셀-연도), 괄호 안은 95% 짝지은 블록 부트스트랩 구간.
\section{결과}\label{sec:results}

\subsection{관측 연도 기후 입력의 효과}
관측은 시간적으로 고르지 않다(그림~\ref{fig:data}b, 표~\ref{tab:data}). 알래스카 셀-연도의 \ntopyrshareAk\%가 \ntopyrAk{}년에 관측되었고, 레나 델타 라벨은 \nnyrLe{}개 해에 분포한다. 측정은 대부분 8월에 이루어졌으며(관측일 중앙값 알래스카 \ndoymedAk, 레나 델타 \ndoymedLe, 캐나다 \ndoymedCa{}번째 날), 관측일까지 누적한 TDD는 계절 전체 값의 중앙값 \ntddfracAk, \ntddfracLe, \ntddfracCa{}배였다(식~(\ref{eq:tddobs})).

평년 기후 입력을 관측 연도 기후 입력으로 바꾸면 레나 델타에서 재보정 Stefan 모델(\ndforcLePo{} cm)과 물리 잔차 모델(\ndforcLeRo{} cm)의, 캐나다에서 물리 잔차 모델(\ndforcCaRo{} cm)의 공간 블록 RMSE가 감소하였다(그림~\ref{fig:accuracy}a). 알래스카에서는 변화가 판정되지 않았다(Stefan \ndforcAkPo{} cm). 라벨이 한 해에 몰려 있어 관측 연도 입력이 공간 분포 이상의 정보를 거의 더하지 못하기 때문으로 해석된다. 관측 연도 입력은 각 예측을 명시된 해의 기후와 연결하므로, 아래에 제시하는 연도별 지도의 전제가 된다.

\paperfig{Fig3_accuracy}{1.0}{\textbf{관측 연도 기후 입력, 검증 설계, 잔차 이득의 조건.} 모든 패널은 2010--2024년 셀-연도 라벨과 겹 밖 예측을 사용한다. \textbf{a}, 평년 기후 입력(2011--2024년 평균)과 관측 연도 기후 입력(관측일까지 누적)의 공간 블록 RMSE. \textbf{b}, 다섯 검증 설계의 RMSE. 지역 홀드아웃에서는 재보정 Stefan 모델이 원천 계수 모델과 같으며, 축 범위 밖 값은 숫자로 표기하였다. \textbf{c}, 공변량 묶음 여섯 가지에서 물리 잔차 모델과 재보정 Stefan 모델의 $\Delta$RMSE. \textbf{d}, 여덟 공변량의 지역 안 3분위에서의 같은 $\Delta$RMSE. 테두리 칸은 95\% 구간이 0을 제외하고, 빗금 칸은 셀-연도 30개 또는 블록 5개 미만이다. 구간은 0.5$^\circ$ 블록 2000회 재추출의 95\% 짝지은 블록 부트스트랩 구간이며, 음수는 물리 잔차 모델이 낫다는 뜻이다.}{fig:accuracy}

\subsection{공간 · 시간 · 지역 검증에서의 정확도}
방법의 순위는 검증 설계에 따라 달라졌다(그림~\ref{fig:accuracy}b, 보충 표 S5). 무작위 교차검증에서 직접 ML의 오차가 가장 낮았고(알래스카 \nrmseAkrandDz{} cm), 공간 블록 교차검증에서는 직접 ML의 오차가 알래스카 \noptimAkDz, 레나 델타 \noptimLeDz, 캐나다 \noptimCaDz{} cm 증가하였다. 연도 홀드아웃에서는 라벨이 여러 해에 분포한 레나 델타에서 직접 ML의 오차가 \nrmseLeblkDz{} cm에서 \nrmseLeyrDz{} cm로 증가하였고, 블록 × 연도 설계는 세 지역 모두에서 지역 라벨로 적합하는 모든 방법의 지역 안 오차가 가장 컸다. 지역 전체를 학습에서 제외하면 레나 델타에서 직접 ML의 오차는 \nrmseLeregDz{} cm로, 다른 지역 계수를 쓴 Stefan 모델의 \nrmseLeregPz{} cm보다 컸다($\Delta$RMSE \ndrefLeregDz{} cm). 같은 시험에서 물리 잔차 모델은 Stefan 모델 수준을 유지하였고(\ndrefLeregRo{} cm), 캐나다에서는 Stefan 모델보다 정확하였다(\ndrefCaregRo{} cm). 따라서 공간 블록 교차검증을 주 설계로, 지역 홀드아웃을 전이성 시험으로 사용한다.

\subsection{물리 잔차 학습의 이득과 그 조건}
관측 연도 입력을 사용한 공간 블록 교차검증에서 물리 잔차 모델은 재보정 Stefan 모델의 RMSE를 레나 델타에서 \nrmseLeblkPo{} cm에서 \nrmseLeblkRo{} cm로($\Delta$RMSE \ndrefLeblkRo{} cm), 캐나다에서 \nrmseCablkPo{} cm에서 \nrmseCablkRo{} cm로(\ndrefCablkRo{} cm) 낮추었고, 알래스카에서는 두 모델 모두 약 \nrmseAkblkRo{} cm였다(\ndrefAkblkRo{} cm). 이득은 잔차 학습기에 주어진 공변량에 따라 달라졌다(그림~\ref{fig:accuracy}c). 기후 공변량만 쓰면 잔차 모델은 알래스카(\ndfeatAkFclim{} cm)와 캐나다(\ndfeatCaFclim{} cm)에서 Stefan 모델을 개선하지 못하였고, 지형과 토양 공변량을 더해야 위의 이득이 나타났다. MODIS 식생 · 눈 소실 시기(2010--2024년)나 고해상 토양 · 피복 분류를 더한 효과는 레나 델타와 캐나다에서 0.5 cm 미만이었고, 알래스카에서는 고해상 분류가 작은 이득을 더하였다(\ndfeatAkFhires{} cm).

결과는 물리 모델이 학습기에 들어가는 방식에 따라 달라졌다(보충 그림 S2). Stefan 유사라벨로 학습 자료를 증강하면 Stefan 모델을 재현하는 데 그쳤고(레나 델타 \ndrefLeblkDo{} cm), Stefan 예측을 입력 특징으로 쓰면 직접 ML과 비슷한 결과를 보였다. 지역의 라벨을 모두 쓰면 직접 ML의 점추정 오차는 세 지역 모두에서 물리 잔차 모델보다 낮았으나(알래스카 \nrmseAkblkDz{} cm 대 \nrmseAkblkRo{} cm, 레나 델타 \nrmseLeblkDz{} cm 대 \nrmseLeblkRo{} cm, 캐나다 \nrmseCablkDz{} cm 대 \nrmseCablkRo{} cm), 재보정한 Stefan 모델과의 차이는 캐나다에서만 판정되었고(\ndrefCablkDz{} cm) 지역 홀드아웃에서는 오차가 커졌다. CatBoost 대신 랜덤 포레스트, 다층 퍼셉트론, TabM, FT-Transformer를 사용하여도 레나 델타와 캐나다에서 물리 잔차 이득의 부호는 바뀌지 않았다.

여덟 공변량의 지역 안 3분위로 나누어 보면, 이득은 레나 델타(\ngcLecells{}개 3분위 가운데 \ngcLesig{}개에서 구간이 0 제외)와 캐나다(\ngcCacells{}개 가운데 \ngcCasig{}개)의 대부분 조건에서 나타났고 Stefan 모델의 오차가 클수록 컸다. 3분위 간 Stefan 오차와 이득의 상관은 레나 델타 \ngcLecorr, 캐나다 \ngcCacorr{}였다(그림~\ref{fig:accuracy}d). 레나 델타의 가장 따뜻한 TDD 3분위에서 이득은 \ndtddLewarm{} cm였으나, 알래스카에서는 어느 3분위도 \ngcAkmaxabs{} cm를 넘게 바뀌지 않았다. 이는 잔차 학습기가 토양과 지형 때문에 실효 Stefan 계수가 지역 값에서 벗어나는 부분을 보정함을 시사한다. 이 해석과 일치하게, 잔차 모델의 평균 절대 SHAP 값 가운데 지형 변수가 레나 델타에서 \nshapLeResTerrain\%, 토양 변수가 캐나다에서 \nshapCaResSoil\%를 차지하였다(보충 그림 S3).

\subsection{표본 효율과 새 지역으로의 전이}
라벨 수 실험은 라벨이 적을 때의 거동을 보여 준다(그림~\ref{fig:labels}a). 한 지역의 2010--2024년 셀-연도에서 학습 라벨 25개를 뽑으면, 물리 잔차 모델은 같은 라벨로 재보정한 Stefan 모델과 같거나 더 정확하였고(알래스카 \nllAkRotwfive, 레나 델타 \nllLeRotwfive, 캐나다 \nllCaRotwfive{} cm), 라벨 100개에서는 세 지역 모두에서 이득의 구간이 0을 제외하였다(\nllAkRohundred, \nllLeRohundred, \nllCaRohundred{} cm). 라벨 25개로 학습한 직접 ML은 알래스카(\nllAkDztwfive{} cm)와 캐나다(\nllCaDztwfive{} cm)에서 Stefan 모델보다 점추정 오차가 컸다. 물리 입력 ML과 물리 유사라벨은 전체 라벨 결과와 같은 경향을 보였다(보충 그림 S4).

\paperfig{Fig4_labels}{1.0}{\textbf{적은 라벨과 새 지역에서의 정확도.} \textbf{a}, 각 지역에서 같은 $n$개 학습 라벨(2010--2024년 라벨, 3회 추출)로 재보정한 Stefan 모델 대비 각 방법의 $\Delta$RMSE. 띠는 물리 잔차 모델의 95\% 블록 부트스트랩 구간이다. \textbf{b}, 다른 두 지역과 대상 지역 라벨 $n$개로 학습한 모델의 지역별 RMSE. \textbf{c}, 외부 CALM 지역(2010--2024년 관측)에서 같은 외부 지역 라벨 수에 따른 RMSE. Wei 등, Liu 등, ESA CCI v5 값이 모두 있는 셀-연도 \nftrows{}개(2010--2020년)에서 계산하였고, 가로선은 제품이며 Wei 등과 Liu 등은 이 지점들을 학습 자료에 포함한다.}{fig:labels}

전이 실험에서는 다른 두 지역과 대상 지역 라벨 $n$개로 학습한 모델로 대상 지역을 예측하였다(그림~\ref{fig:labels}b). 대상 라벨이 없을 때 레나 델타에서 직접 ML은 \ntlrmseLeDszero{} cm로, 원천 계수 Stefan 모델(\ntlrmseLePzzero{} cm)과 물리 잔차 모델(\ntlrmseLeRszero{} cm)보다 오차가 컸다. 대상 라벨 25개를 더하면 물리 잔차 모델은 \ntlrmseLeRstwfive{} cm, 직접 ML은 \ntlrmseLeDstwfive{} cm로 감소하였다. 알래스카에서는 물리 잔차 모델이 모든 $n$에서 가장 정확하였다(대상 라벨 0개에서 \ntlrmseAkRszero{} cm). 캐나다에서는 원천 지역으로 학습한 직접 ML이 가장 정확하였고(\ntlrmseCaDszero{} cm), 캐나다 라벨 25개로 Stefan 계수를 재보정하면 오차가 \ntlrmseCaPzzero{} cm에서 \ntlrmseCaPntwfive{} cm로 증가하였으며 이 기준값을 쓰는 물리 잔차 모델도 \ntlrmseCaRszero{} cm에서 \ntlrmseCaRstwfive{} cm로 증가하였다. 캐나다 지점들의 융해 깊이 범위가 넓기 때문이다.

같은 절차를 모델 개발에 쓰지 않은 외부 CALM 지역 8곳의 2010--2024년 채점 셀-연도 \nftScoredRows{}개에 적용하였다(그림~\ref{fig:labels}c, 보충 그림 S5). Wei 등, Liu 등, ESA CCI v5 값이 모두 있는 7개 지역의 셀-연도 \nftrows{}개에서, 세 연구 지역으로 학습한 물리 잔차 모델은 지역 라벨 없이 \nftAllRszero{} cm였다. 같은 외부 지역의 라벨 5개를 더하면 \nftAllRsfive{} cm, 25개를 더하면 \nftAllRstwfive{} cm로 감소하였고, 직접 ML은 라벨 25개에서 \nftAllDstwfive{} cm였다. 라벨 5개에서 물리 잔차 모델은 ALT 관측으로 학습하지 않은 ESA CCI v5보다 정확하였다(같은 \nftrows{}개 셀-연도에서 \nftAllCcifivezero{} cm, ESA CCI v5 값이 있는 셀-연도 \nftdCciRows{}개 전체에서 $\Delta$RMSE \nftdRsCcifivefive{} cm). 이 지점들을 학습 자료에 포함한 Wei 등과 Liu 등은 \nftAllWeizero{} cm와 \nftAllLiuzero{} cm였다. 오차 감소의 대부분은 Stefan 계수의 재보정에서 나왔고, 잔차 학습기는 이 지점들에서 오차를 추가로 줄이지 못하였다(셀-연도 \nftScoredRows{}개 전체, 라벨 5개에서 \nftdRsPnfive{} cm). 합산 값에는 지역별로 다른 반응이 섞여 있다. 라벨 5개에서 오차는 서시베리아, 중앙 시베리아, 유럽 러시아, 스발바르에서 줄었으나, 라벨 없이도 오차가 10 cm 미만이던 추코트카(\nftChRszero{} cm에서 \nftChRsfive{} cm)와 캄차카(\nftKaRszero{} cm에서 \nftKaRsfive{} cm), 그리고 북동 시베리아(\nftNsRszero{} cm에서 \nftNsRsfive{} cm)에서는 늘었다. 서 · 중부 시베리아와 스발바르의 활동층이 연구 지역의 범위보다 두껍기 때문이다.

\subsection{같은 셀 · 같은 해의 공개 제품 비교}
공통 셀-연도(\nyrscc{}년)에서 물리 잔차 모델은 레나 델타에서 \nccLeRoall{} cm로, ESA CCI v5(\nccLeCcifiveall{} cm), Aalto 등(\nccLeAaltoall{} cm), ESA CCI v4(\nccLeCcifourall{} cm), Liu 등(\nccLeLiuall{} cm), Wei 등(\nccLeWeiall{} cm)보다 정확하였다(그림~\ref{fig:products}a, 보충 표 S6). 제품 학습 지점에서 25 km 안의 셀을 제외하면 짝지은 차이는 Wei 등 대비 \ndprodLeweiRomask{} cm, ESA CCI v5 대비 \ndprodLeccifiveRomask{} cm였다. 알래스카(\nccAkRoall{} cm 대 Wei 등 \nccAkWeiall{} cm, Liu 등 \nccAkLiuall{} cm)와 캐나다(\nccCaRoall{} cm 대 \nccCaWeiall{} cm, \nccCaLiuall{} cm)에서는 모델과 두 ML 제품이 같은 수준이었다.

\paperfig{Fig5_products}{1.0}{\textbf{같은 셀 · 같은 해의 공개 제품 비교.} \textbf{a}, 공통 셀-연도(\nyrscc{}년)에서 세 방법(공간 블록 겹 밖 예측)과 제품 5종의 RMSE. 모든 셀과 Wei 등 학습 지점에서 25 km 넘게 떨어진 셀을 구분하였다. ESA CCI v5와 Aalto 등은 값이 있는 셀-연도에서 채점하였고, 축 밖 값은 숫자로 적었다. \textbf{b}, 해마다 같은 셀-연도에서 계산한 관측 연도별 RMSE. 셀-연도 30개 이상인 해(알래스카 \nyrsywAkko, 레나 델타 \nyrsywLeko, 캐나다 \nyrsywCako)만 나타내었고, Liu 등은 2020년까지이며, 막대는 해마다의 셀-연도 수이다.}{fig:products}

이러한 지역 차이는 제품의 학습 자료와 대응한다. 알래스카 라벨 셀은 Wei 등의 학습 지점에서 중앙값 \ndistweiAkmed{} km 떨어져 있고 Liu 등은 같은 알래스카 · 캐나다 현장 자료의 이전 판으로 학습한 반면, 레나 델타 셀은 중앙값 \ndistweiLemed{} km 떨어져 있다(보충 그림 S6). 두 제품 모두 델타 북부의 얇은 활동층을 과대 추정하며, 북위 72$^\circ$ 이북의 평균 편향은 \nbiasLeNorthWei{} cm와 \nbiasLeNorthLiu{} cm이다. 레나 델타 라벨 없이 학습한 모델도 같은 셀-연도에서 이미 Wei 등보다 정확하였으므로(\ntlvprmseLeRsWeizero{} cm 대 \ntlvprmseLeWeizero{} cm, $\Delta$RMSE \ntlvpLeRsWeizero{} cm), 이 지역의 우위는 지역 라벨에만 의존하지 않는다.

같은 해의 지도는 이 차이를 공간적으로 보여 준다(그림~\ref{fig:prodmaps}). 2014년 알래스카에서는 세 지도가 관측 셀에서 비슷하였다(RMSE: 우리 모델의 겹 밖 예측 \nfeAkOof{} cm, Wei 등 \nfeAkWei{} cm, Liu 등 \nfeAkLiu{} cm). 이 관측 셀 \nfeAkN{}개 중 \nfeAkNearWei{}개는 Wei 등의 학습 지점에서 25 km 안에 있다. 지도 간 차이는 적용 가능 영역 밖인 고도 600 m 이상의 산지에서 가장 컸다. 북위 68--70$^\circ$에서 제품은 우리 지도보다 \nmapAkFootMtnWei{} cm(Wei 등), \nmapAkFootMtnLiu{} cm(Liu 등) 깊었다. 관측 셀 \nfeLeN{}개가 모두 이 학습 지점에서 25 km보다 먼 2016년 레나 델타에서는 두 제품이 델타 전반에서 관측보다 깊었고(평균 편차 \nfeLeBiasWei{} cm, \nfeLeBiasLiu{} cm, RMSE \nfeLeWei{} cm, \nfeLeLiu{} cm), 우리 모델의 RMSE는 \nfeLeOof{} cm였다.

\paperfig{Fig6_product_maps}{1.0}{\textbf{같은 해의 물리 잔차 모델 지도와 공개 제품 지도.} \textbf{a--c}, 2014년 알래스카의 물리 잔차 모델, Wei 등, Liu 등 지도(같은 색 범위). \textbf{d}, 2014년 관측 셀 중 세 지도 모두 값이 있는 \nfeAkN{}개 셀의 예측 빼기 관측 ALT와 지도별 RMSE. 우리 모델의 예측은 공간 블록 겹 밖 값이다. 축 범위 밖 값은 숫자로 표기하였다. 제품 값에는 학습 지점 마스크를 적용하지 않았으며, \nfeAkN{}개 중 \nfeAkNearWei{}개 셀은 Wei 등의 학습 지점에서 25 km 안에 있다. \textbf{e--h}, 2016년 레나 델타의 같은 비교(\nfeLeN{}개 셀, 모두 학습 지점에서 25 km보다 멂). 지도 범위는 ESA CCI 영구동토 확률 50\% 이상인 지역이며, 1 km 값을 표시용으로 250 m에 보간하였다.}{fig:prodmaps}

연도별로 보면 물리 잔차 모델은 셀-연도 30개 이상인 레나 델타의 \nywLen{}개 해(\nyrsywLeko) 모두에서 Wei 등과 Stefan 모델보다 오차가 낮았고, 알래스카(\nyrsywAkko)에서는 Wei 등 대비 \nywAkRoWein{}개 해 가운데 \nywAkRoWeiwin{}개 해, Stefan 모델 대비 \nywAkRoPon{}개 해 가운데 \nywAkRoPowin{}개 해, 캐나다에서는 \nywCaRoPon{}개 해 가운데 Stefan 모델 대비 \nywCaRoPowin{}개 해에서 낮았다(그림~\ref{fig:products}b). 여러 해 관측된 알래스카 셀에서 관측과 예측의 연도 간 변화 상관은 물리 잔차 모델 \nrdeltaAkRo, Wei 등 \nrdeltaAkWei, Liu 등 \nrdeltaAkLiu{}였으며, 레나 델타에서는 제품이 연도 간 변화를 더 잘 따랐다(\nrdeltaLeWei{}와 \nrdeltaLeLiu{} 대 \nrdeltaLeRo).

잔차 학습기는 제품을 기준값으로 하여 적용할 수도 있다(식~(\ref{eq:prod})). 레나 델타 라벨로 Wei 등의 재보정 제품에 잔차 학습기를 더하면 오차가 \npadLeRpPcal{} cm 변하였고, Stefan 기준 모델은 제품 기준 모델보다 여전히 정확하였다(\npadLeRoRp{} cm, 보충 그림 S7). 알래스카와 캐나다에서는 두 기준값의 차이가 구분되지 않았다.

\subsection{연도가 명시된 1 km ALT 지도(2011--2024년)}
알래스카와 레나 델타의 모든 셀-연도로 학습한 물리 잔차 모델을 2011--2024년 각 해의 ERA5-Land 기후로 적용하였다(그림~\ref{fig:maps}). 지도는 알래스카 \nncellmapAk{}개, 레나 델타 \nncellmapLe{}개 셀이다. 2011--2024년 평균 대비 지역 평균 편차는 알래스카에서 \nanminyrAk{}년 \nanminAk{} cm부터 \nanmaxyrAk{}년 \nanmaxAk{} cm까지, 레나 델타에서 \nanminyrLe{}년 \nanminLe{} cm부터 \nanmaxyrLe{}년 \nanmaxLe{} cm까지였으며, 공간 분포도 해마다 달랐다(그림~\ref{fig:maps}b). 레나 델타에서 관측 셀의 평균 편차는 셀-연도 50개 이상인 \nanobsnLe{}개 해(\nanobsyrsLe{}년) 가운데 \nanobssameLe{}개 해에서 예측 지역 편차와 부호가 같았다($r$ = \nanobsrLe). 관측 편차가 \nanobsLeTwentyObs{} cm(셀-연도 \nanobsLeTwentyN{}개)에 이른 2020년을 빼도 상관은 \nanobsrLeNoTwenty{}였다. 해마다 관측 셀 구성이 바뀌는 알래스카에서는 일치가 약하였고($r$ = \nanobsrAk), 예측의 연도 간 범위(\nanpredrangeAk{} cm)는 관측 셀 평균의 범위(\nanobsrangeAk{} cm)보다 작았다(그림~\ref{fig:maps}c). 적용 가능 영역은 알래스카 셀의 \naoaAk\%, 레나 델타 셀의 \naoaLe\%를 포함하며, 영역 밖의 셀은 주로 산지와 영구동토 지대의 남쪽 가장자리에 있다(보충 그림 S8). 90\% 예측 구간의 반폭은 \nhwAk{} cm와 \nhwLe{} cm이다. 잔차 학습기가 더하는 보정량은 대부분 $\pm$10 cm 안에 있고(보충 그림 S9), 장기 관측 셀의 시계열은 보충 그림 S10에 나타내었다.

\paperfig{Fig7_maps}{1.0}{\textbf{연도가 명시된 1 km ALT 지도(2011--2024년).} \textbf{a}, 연도별 지도로 계산한 알래스카 물리 잔차 모델의 2011--2024년 평균 ALT. 각 연도 지도는 그 해의 계절 전체 ERA5-Land 기후를 사용하며, 옅은 색은 적용 가능 영역 밖, 회색 육지는 지도 범위 밖이다. \textbf{b}, 2011--2024년 평균 대비 각 해의 ALT 편차(모든 해 같은 색 범위 $\pm$15 cm). \textbf{c}, 예측 지도의 지역 평균 편차와 셀-연도 50개 이상인 해의 관측 셀 평균 편차(알래스카 \nanobsyrsAk{}년, 레나 델타 \nanobsyrsLe{}년). 관측 편차는 각 관측 셀의 2011--2024년 지도 평균을 기준으로 계산하고 같은 해들의 예측 편차 평균에 중심을 맞추었으며, 레나 델타 2020년 값에는 셀-연도 수를 표기하였다. 1 km 값은 표시용으로 250 m에 보간하였다.}{fig:maps}
